% =============================================================
%  Appendix A1 -- Signal-Extraction Model Derivation
%  Label: app:model
% =============================================================

\section{Signal-Extraction Model Derivation}\label{app:model}

This appendix provides the full derivation of the rational
benchmark and diagnostic extension summarised in
Section~\ref{sec:framework}.

\subsection{Environment}

A household in province $p$ at wave $t$ observes a meat-price
signal $s_{pt}$ and must forecast headline CPI inflation
$\pi_{p,\,t \to t+1}$ over the next inter-wave period.
The signal $s_{pt}$ is the cumulative log change in the
province-level meat CPI between waves $t-1$ and $t$, as
defined in equation~\eqref{eq:meat_shock_data}.

Headline CPI responds to the meat shock according to
\[
  \pi_{p,\,t \to t+1}
  = \bar{\pi} + \omega\,\rho\,s_{pt} + \eta_{pt},
\]
where $\omega \in (0,1)$ is the CPI weight on meat,
$\rho \in [0,1]$ is the persistence of the shock
(the fraction that carries forward into the next period),
and $\eta_{pt} \sim (0, \sigma_\eta^2)$ is an aggregate
noise term orthogonal to the meat signal.
The noise term $\eta_{pt}$ captures all non-meat determinants
of headline inflation: monetary policy, non-food supply shocks,
demand conditions, and measurement error in the CPI.

The assumption that headline CPI is linear in the meat shock is
a first-order approximation.
In the CPI accounting identity, the contribution of meat to
headline inflation is $\omega \times \pi^{\text{meat}}$, where
$\pi^{\text{meat}}$ is meat-price inflation.
If the current meat shock $s_{pt}$ persists at rate $\rho$
into the next period, the predicted contribution is
$\omega \cdot \rho \cdot s_{pt}$.
The linearity assumption also rules out threshold effects
(e.g., overreaction only when meat-price changes exceed
some attention threshold) and interaction effects between
meat prices and other CPI components.
These features could be added at the cost of additional
parameters that the data cannot identify separately.

The household does not know $\rho$ with certainty.
It holds a point belief $\hat{\rho} \in [0,1]$ about
persistence.
This parametrises the household's uncertainty about whether
the current meat-price movement will persist or reverse.
The household knows $\omega$ and observes $s_{pt}$ without
error.
Relaxing the assumption that $\omega$ is known shifts the
rational bound proportionally, as discussed in
Section~\ref{sec:transitory}.

\subsection{Rational Forecast}

Under rational expectations given the household's belief
about persistence, the optimal forecast is
\[
  \E_t^R[\pi_{p,\,t \to t+1}]
  = \bar{\pi} + \omega\,\hat{\rho}\,s_{pt}.
\]
This is the minimum-mean-squared-error forecast given the
information set $\{s_{pt}, \omega, \hat{\rho}\}$ and the
linear signal structure.
The household weights the observed meat shock by its CPI
weight and its perceived persistence.

The forecast error is
\begin{align}
  \mathrm{FE}^R_{pt}
  &= \pi_{p,\,t \to t+1}
    - \E_t^R[\pi_{p,\,t \to t+1}] \notag \\
  &= \omega(\rho - \hat{\rho})\,s_{pt} + \eta_{pt}.
    \label{eq:app_rational_fe}
\end{align}
The forecast error is mean-zero conditional on $s_{pt}$ only
if $\hat{\rho} = \rho$---that is, only if the household has
the correct belief about persistence.
If $\hat{\rho} \neq \rho$, the forecast error covaries with
the signal, and the regression coefficient
$\beta_3^R = \omega(\rho - \hat{\rho})$ is nonzero.

The absolute value of the rational coefficient satisfies
\[
  |\beta_3^R| = \omega\,|\rho - \hat{\rho}|
  \leq \omega
  \approx 0.05,
\]
since $\rho, \hat{\rho} \in [0,1]$.
The maximum is attained at the corners: when the household
believes the shock is permanent ($\hat{\rho} = 1$) but it is
fully transitory ($\rho = 0$), or vice versa.
This bound is the sharpest test the model delivers.
It does not require knowledge of $\hat{\rho}$: regardless of
what the household believes, the rational forecast-error
coefficient is bounded by $\omega$.

\subsection{Diagnostic Forecast}

Under diagnostic expectations \citep{Bordalo2018}, the
household distorts its likelihood by overweighting states made
more representative by the recent signal.
In the Gaussian case, this amounts to an additive distortion
with parameter $\theta > 0$:
\[
  \E_t^\theta[\pi_{p,\,t \to t+1}]
  = \bar{\pi} + \omega\,\hat{\rho}\,s_{pt}
    + \theta\,\omega\,s_{pt}.
\]
The diagnostic distortion is proportional to the signal $s_{pt}$
and to the CPI weight $\omega$.
The parameter $\theta$ governs the degree of overweighting:
$\theta = 0$ reduces to the rational case, and $\theta > 0$
generates over-reaction in the direction of the signal.

The diagnostic forecast error is
\begin{align}
  \mathrm{FE}^\theta_{pt}
  &= \omega(\rho - \hat{\rho})\,s_{pt}
    - \theta\,\omega\,s_{pt} + \eta_{pt} \notag \\
  &= \omega(\rho - \hat{\rho} - \theta)\,s_{pt}
    + \eta_{pt}.
    \label{eq:app_diagnostic_fe}
\end{align}
The regression coefficient is
$\beta_3^\theta = \omega(\rho - \hat{\rho} - \theta)$.
Since $\theta > 0$:
\[
  |\beta_3^\theta| > |\beta_3^R|
  \quad \text{for all } \rho, \hat{\rho} \in [0,1].
\]
The gap $|\beta_3^\theta| - |\beta_3^R| = \omega\,\theta$
identifies the diagnostic parameter scaled by the CPI weight.

\subsection{Time-Aggregation Extension}

The one-period model treats the meat shock as a single
realisation.
In the data, the treatment $S_{pt} = \sum_{m=1}^{M} s_{pm}$
is a cumulative shock over $M \approx 24$ months, and the
forecast horizon covers $H \approx 24$ forward months.
This subsection derives the rational bound for the
time-aggregated object.

Under the AR(1) specification, the headline CPI contribution
of a monthly meat shock $s_{pm}$ at horizon $h$ months ahead
is $\omega\,\rho^h\,s_{pm}$.
The total forward headline CPI response to the sequence of
monthly shocks is
\[
  \sum_{m=1}^{M}\sum_{h=1}^{H}
  \omega\,\rho^h\,s_{pm}
  = \omega\,\frac{\rho(1-\rho^H)}{1-\rho}\,S_{pt},
\]
where I use the fact that
$\sum_{h=1}^H \rho^h = \rho(1-\rho^H)/(1-\rho)$ and
$\sum_{m=1}^M s_{pm} = S_{pt}$.

A rational household with belief $\hat{\rho}$ forecasts the
same expression with $\hat{\rho}$ replacing $\rho$.
The forecast-error regression coefficient on $S_{pt}$ is
\[
  \beta_3^R
  = \omega\left[
      \frac{\rho(1-\rho^H)}{1-\rho}
    - \frac{\hat{\rho}(1-\hat{\rho}^H)}{1-\hat{\rho}}
    \right].
\]

Define $g(\rho) \equiv \rho(1-\rho^H)/(1-\rho)$.
For $\rho \in [0,1]$ and finite $H$, $g$ is continuous with
$g(0) = 0$ and $\lim_{\rho \to 1} g(\rho) = H$.
The range of $g$ is $[0, H]$.
Therefore
$|g(\rho) - g(\hat{\rho})| \leq H$ for all
$\rho, \hat{\rho} \in [0,1]$, and
\[
  |\beta_3^R| \leq \omega\,H.
\]
With $H \approx 24$, the time-aggregated rational bound is
$24 \times 0.05 = 1.20$.

The bound is tightest when the treatment and forecast horizons
are equal ($M \approx H$), which is the case in the CFPS
design.
If the treatment were concentrated in a single month ($M = 1$)
but the forecast horizon remained 24 months, the bound per unit
of $S_{pt}$ would be $\omega H = 1.20$, unchanged.
The estimated $|\hat{\beta}_3| = 7.13$ exceeds this bound by a
factor of six.

\subsection{Calibration}

I estimate $\rho$ from a pooled AR(1) regression on monthly
province-level meat CPI log changes with province fixed effects.
The estimating equation is
\[
  \Delta \log(\mathrm{MeatCPI}_{pm})
  = \alpha_p
  + \rho\,\Delta \log(\mathrm{MeatCPI}_{p,m-1})
  + \varepsilon_{pm},
\]
with standard errors clustered at the province level.
The sample covers January 2016 through December 2025 for 31
provinces, yielding approximately 3,689 monthly observations.
The estimated persistence is $\hat{\rho} = 0.480$.

Given $\hat{\rho}$ from the data and $\omega = 0.05$, the
rational benchmark is
$|\beta_3^R| \leq \omega = 0.05$.
The estimated $\hat{\beta}_3 = -7.13$ under the preferred
specification yields an implied diagnostic parameter that is
large.
Even using the time-aggregated rational bound of $1.20$ as
the baseline, the excess is
$(7.13 - 1.20)/0.05 = 119$.
Because the CFPS expectation is ordinal, the exact magnitude
of $\hat{\beta}_3$ embeds an arbitrary scaling choice, and
the implied $\hat{\theta}$ should not be taken at face value.
The qualitative conclusion is that $\hat{\theta} \gg 0$:
overreaction is large relative to the rational benchmark.
A scaling-free lower bound on the excess comes from the
pass-through coefficient: $|\hat{\beta}_2| = 5.25$ exceeds
the time-aggregated rational bound of 1.20 without involving
any ordinal variable (Appendix~\ref{app:scaling}).

\subsection{Comparative Statics}

The gap between the diagnostic and rational forecast errors is
$|\beta_3^\theta| - |\beta_3^R| = \omega\,\theta$.
This expression shows how the observable wedge depends on the
model parameters.

The wedge is increasing in $\omega$: a higher CPI weight on
meat amplifies the overreaction because each unit of the
diagnostic distortion translates into a larger revision of
the expected headline inflation rate.
In a country where meat carries a CPI weight of 0.10 rather
than 0.05, the same $\theta$ would produce a wedge twice as
large.
This has implications for external validity: the overreaction
documented here should be larger in economies with high food
CPI weights and smaller in economies where food is a minor
CPI component.

The wedge is increasing in $\theta$: a more diagnostic
household produces a larger forecast error for a given shock.
The $\theta$ parameter is treated as a constant in the
baseline model, but it could vary across households (for
example, with education or financial literacy) or across
signal types (for example, more salient signals could
generate a higher effective $\theta$).
The null heterogeneity results in Table~\ref{tab:heterogeneity}
suggest that $\theta$ does not vary strongly with observable
demographics in this sample, though this does not rule out
heterogeneity along unobserved dimensions such as news
consumption or market experience.

The wedge does not depend on $\rho$ or $\hat{\rho}$.
The persistence parameters affect both the rational and
diagnostic coefficients, but their effects cancel in the
difference.
This is a useful property: the identification of $\theta$
does not require the econometrician to know or estimate the
household's belief about persistence.
In practice, the calibration exercise in
Section~\ref{sec:calibration} uses $\rho$ to compute the
rational bound separately, but the excess beyond the bound
identifies $\omega\,\theta$ regardless of $\rho$.

\subsection{Connection to the Empirical Specification}

The model maps directly to the regression equations in
Section~\ref{sec:three_equations}.
The meat shock $s_{pt}$ corresponds to the
$\mathrm{MeatShock}_{p,\,t-1\to t}$ variable constructed from
NBS data.
The forecast error $\mathrm{FE}_{pt}$ corresponds to
$\pi^{\text{headline}}_{p,\,t\to t+1} - \tilde{\mu}_{ipt}$
in equation~\eqref{eq:eq3}.

The model predicts that a regression of the forecast error
on the meat shock should yield a coefficient $\beta_3$ that
is bounded by $\omega$ under rationality.
The empirical equation~\eqref{eq:eq3} includes fixed effects
$\delta_{r(p)\times t}$ that are not present in the model.
In the model, $\bar{\pi}$ is a constant; in the data, the
fixed effects absorb variation in $\bar{\pi}$ across regions
and waves, as well as any region-wave-specific component of
$\eta_{pt}$.
The fixed effects therefore improve the precision of
$\hat{\beta}_3$ without affecting its probability limit,
provided that the meat shock is conditionally uncorrelated
with the fixed-effect-absorbed component of $\eta_{pt}$.

The household controls $X_{ipt}$ in the empirical specification
have no counterpart in the model.
They are included to absorb observable heterogeneity in
expectation formation (e.g., education, income, urban status)
that might be correlated with provincial meat-shock exposure.
Their inclusion does not alter the theoretical prediction;
it sharpens the empirical estimate by reducing residual
variance.

One gap between model and data deserves note.
The model assumes that the household observes a continuous
signal $s_{pt}$ and forms a continuous expectation.
The CFPS expectation is ordinal: $\mu_{ipt} \in \{-1,0,1\}$.
No information in the survey resolves the mapping from
ordinal codes to CPI units.
The overreaction test does not require this mapping:
the sign of $\hat{\beta}_3$ is invariant to any monotone
rescaling of the ordinal variable
(Appendix~\ref{app:scaling}), and the pass-through coefficient
$\hat{\beta}_2$ provides a scaling-free magnitude test.
The implied $\hat{\theta}$ depends on the assumed mapping
and should not be interpreted as a point estimate.

\subsection{Limitations of the Model}

The signal-extraction framework is deliberately parsimonious.
Its purpose is to provide a quantitative benchmark for the
rational forecast error, not to describe the full
information environment of Chinese households.
Several features are omitted.

The model assumes a single signal.
In reality, households observe multiple price signals (meat,
grain, vegetables, housing, services) and must aggregate them
into an inflation forecast.
The prediction that the forecast-error coefficient on the meat
shock is bounded by $\omega$ holds in a multi-signal extension
provided that the meat shock is orthogonal to other signals
conditional on fixed effects; the fixed-effect structure in
the empirical specification is designed to absorb common
components across signals.

The model also assumes static expectations: the household
forms a one-period-ahead forecast and does not revise it.
In practice, households may update their beliefs between
survey waves as new information arrives.
The biennial CFPS structure does not allow observation of
intra-wave updating, and a dynamic extension with sequential
signals would add parameters the current data cannot identify.

Finally, the model does not endogenise attention.
The diagnostic parameter $\theta$ is a constant that applies
to the meat signal regardless of its size.
A richer model in the spirit of rational inattention
\citep{Sims2003} would allow $\theta$ to depend on the
signal-to-noise ratio or the cost of processing information.
Such an extension could explain why overreaction is larger
during the ASF episode (when meat-price signals were
unusually large) than in normal times, but it requires
additional structure and cross-sectional variation in
attention costs that are not available in the current data.

The model also assumes homogeneous households.
In the data, households differ in income, education, and
geographic location.
A heterogeneous-agent extension would allow $\theta$ or
$\hat{\rho}$ to vary with household characteristics.
The null heterogeneity results suggest that any such
variation is small, but the ordinal expectation measure may
lack the resolution to detect moderate heterogeneity in the
diagnostic parameter.
